\section{Рабочий проект}

\subsection{Спецификация модулей программной системы}

В данном разделе приведено описание основных программных модулей, разработанных в процессе реализации web-сервиса для размещения объявлений об услугах. Программная система построена как PHP-приложение без использования фреймворка. Серверная логика распределена по функциональным каталогам, каждый из которых отвечает за отдельную группу пользовательских сценариев: регистрацию и авторизацию, управление профилем, работу с услугами, настройку расписания, бронирование, переписку, избранное и отзывы.

Клиентская часть представлена HTML-страницами, CSS-стилями и JavaScript-кодом. Хранение данных осуществляется в базе данных MySQL, доступ к которой выполняется через PDO. Взаимодействие пользователя с системой происходит через web-интерфейс и HTTP-запросы.

В проекте используются два способа организации серверной логики. Часть повторно используемых действий вынесена в отдельные PHP-функции, расположенные в общих файлах проекта. Основные пользовательские сценарии реализованы процедурным кодом в отдельных PHP-скриптах-обработчиках. Поэтому в данном разделе отдельно приведены таблицы обычных PHP-функций и процедурных функций модулей.

Под процедурной функцией понимается логически завершённая операция, выполняемая соответствующим скриптом в рамках пользовательского сценария. В таблицах обычных PHP-функций указывается область применения функции, так как данные функции объявлены вне классов и не имеют модификаторов доступа. В таблицах процедурных функций указывается уровень пользовательского доступа к операции: public -- доступно через пользовательский интерфейс, customer -- доступно только заказчику, provider -- доступно только исполнителю, internal -- выполняется внутри серверного сценария и не вызывается пользователем напрямую.

\subsubsection{Модуль общих настроек и вспомогательных функций}

Модуль общих настроек и вспомогательных функций обеспечивает подготовку окружения для работы остальных частей системы. Он включает подключение к базе данных, запуск пользовательской сессии, проверку авторизации, перенаправление пользователей, безопасный вывод данных и валидацию основных пользовательских полей.

Обычные PHP-функции модуля приведены в таблице \ref{common-php-functions:table}.

\begin{xltabular}{\textwidth}{|p{0.3\textwidth}|p{0.22\textwidth}|X|}
	\caption{Обычные PHP-функции модуля общих настроек и вспомогательных функций}
	\label{common-php-functions:table}\\
	\hline
	Функция & Область применения & Описание \\
	\hline
	\endfirsthead
	\caption*{Продолжение таблицы \ref{common-php-functions:table}}\\
	\hline
	Функция & Область применения & Описание \\
	\hline
	\endhead
	\texttt{redirect(\$url)} & Навигация & Выполняет перенаправление пользователя на указанную страницу и завершает выполнение текущего скрипта. \\
	\hline
	\texttt{isLoggedIn()} & Авторизация & Проверяет наличие данных авторизованного пользователя в сессии. \\
	\hline
	\texttt{getUserId()} & Авторизация & Возвращает идентификатор текущего пользователя из данных сессии. \\
	\hline
	\texttt{requireAuth()} & Контроль доступа & Запрещает доступ неавторизованным пользователям к защищённым разделам системы. \\
	\hline
	\texttt{requireGuest()} & Контроль доступа & Ограничивает доступ авторизованных пользователей к страницам регистрации и входа. \\
	\hline
	\texttt{getCurrentUser(\$pdo)} & Профиль пользователя & Получает из базы данных актуальные сведения о текущем пользователе, включая данные профиля и город. \\
	\hline
	\texttt{e(\$string)} & Безопасность вывода & Выполняет экранирование пользовательских данных перед выводом на страницу. \\
	\hline
	\texttt{validateFullName()} & Валидация данных & Проверяет корректность введённого ФИО пользователя. \\
	\hline
	\texttt{validateEmailAddress()} & Валидация данных & Проверяет корректность адреса электронной почты. \\
	\hline
	\texttt{validatePhoneNumber()} & Валидация данных & Проверяет соответствие номера телефона допустимому формату. \\
	\hline
	\texttt{validatePasswordValue()} & Валидация данных & Проверяет длину и сложность пароля пользователя. \\
	\hline
\end{xltabular}

Процедурные функции модуля приведены в таблице \ref{common-procedural-functions:table}.

\begin{xltabular}{\textwidth}{|p{0.3\textwidth}|p{0.2\textwidth}|X|}
	\caption{Процедурные функции модуля общих настроек и вспомогательных функций}
	\label{common-procedural-functions:table}\\
	\hline
	Функция & Уровень\newline доступа & Описание \\
	\hline
	\endfirsthead
	\caption*{Продолжение таблицы \ref{common-procedural-functions:table}}\\
	\hline
	Функция & Уровень\newline доступа & Описание \\
	\hline
	\endhead
	Подключение к базе данных & internal & Формирует параметры соединения с MySQL, создаёт объект PDO, задаёт режим обработки ошибок и формат выборки данных. \\
	\hline
	Запуск пользовательской сессии & internal & Инициализирует механизм сессий PHP, необходимый для хранения данных авторизованного пользователя и временных сообщений. \\
	\hline
	Формирование общей шапки сайта & public & Подключает стили, отображает навигационное меню, ссылки на основные разделы и элементы интерфейса с учётом роли пользователя. \\
	\hline
	Формирование нижней части сайта & public & Завершает HTML-структуру страницы и подключает общие JavaScript-файлы. \\
	\hline
\end{xltabular}

\subsubsection{Модуль регистрации и авторизации пользователей}

Модуль регистрации и авторизации предназначен для создания учётных записей, входа пользователей в систему и завершения сеанса работы. При регистрации пользователь указывает персональные данные, город, роль и пароль. Пароль сохраняется в базе данных в виде хэша, а при входе проверяется средствами PHP.

В данном модуле основные операции реализованы процедурным кодом в скриптах регистрации, входа и выхода. Процедурные функции модуля приведены в таблице \ref{auth-procedural-functions:table}.

\begin{xltabular}{\textwidth}{|p{0.3\textwidth}|p{0.2\textwidth}|X|}
	\caption{Процедурные функции модуля регистрации и авторизации}
	\label{auth-procedural-functions:table}\\
	\hline
	Функция & Уровень\newline доступа & Описание \\
	\hline
	\endfirsthead
	\caption*{Продолжение таблицы \ref{auth-procedural-functions:table}}\\
	\hline
	Функция & Уровень\newline доступа & Описание \\
	\hline
	\endhead
	Отображение формы регистрации & public & Выводит форму создания учётной записи, список городов и варианты выбора роли пользователя. \\
	\hline
	Проверка регистрационных данных & internal & Проверяет ФИО, email, телефон, город, роль, пароль и подтверждение пароля. \\
	\hline
	Проверка уникальности email & internal & Выполняет поиск пользователя с указанным email и предотвращает повторную регистрацию. \\
	\hline
	Создание учётной записи & internal & Хэширует пароль и сохраняет данные нового пользователя в базе данных. \\
	\hline
	Отображение формы входа & public & Выводит форму авторизации пользователя по email и паролю. \\
	\hline
	Проверка данных входа & internal & Находит пользователя по email и проверяет введённый пароль с сохранённым хэшем. \\
	\hline
	Создание пользовательской сессии & internal & Сохраняет в сессии идентификатор пользователя, ФИО, email и роль. \\
	\hline
	Завершение сеанса & public & Удаляет данные пользователя из сессии и перенаправляет его на главную страницу. \\
	\hline
\end{xltabular}

\subsubsection{Модуль управления профилем пользователя}

Модуль управления профилем предназначен для отображения и редактирования персональных данных пользователя. В профиле хранятся ФИО, email, телефон, город, описание и изображение пользователя.

Основные операции модуля реализованы процедурным кодом. Процедурные функции модуля приведены в таблице \ref{profile-procedural-functions:table}.

\begin{xltabular}{\textwidth}{|p{0.3\textwidth}|p{0.2\textwidth}|X|}
	\caption{Процедурные функции модуля управления профилем пользователя}
	\label{profile-procedural-functions:table}\\
	\hline
	Функция & Уровень\newline доступа & Описание \\
	\hline
	\endfirsthead
	\caption*{Продолжение таблицы \ref{profile-procedural-functions:table}}\\
	\hline
	Функция & Уровень\newline доступа & Описание \\
	\hline
	\endhead
	Отображение личного кабинета & public & Выводит сведения о текущем пользователе и доступные действия с учётом его роли. \\
	\hline
	Отображение формы редактирования & public & Предоставляет пользователю форму изменения профиля и список доступных городов. \\
	\hline
	Проверка данных профиля & internal & Проверяет корректность ФИО, телефона, города и текстового описания. \\
	\hline
	Загрузка изображения профиля & internal & Проверяет тип и размер изображения, сохраняет файл в каталог загрузок и обновляет путь к нему в базе данных. \\
	\hline
	Обновление профиля & internal & Сохраняет изменённые данные пользователя в базе данных и обновляет сведения в текущей сессии. \\
	\hline
\end{xltabular}

\subsubsection{Модуль каталога и просмотра услуг}

Модуль каталога и просмотра услуг отвечает за отображение опубликованных услуг, применение поисковых фильтров, просмотр карточки услуги и вывод свободных временных слотов.

Процедурные функции модуля приведены в таблице \ref{catalog-procedural-functions:table}.

\begin{xltabular}{\textwidth}{|p{0.3\textwidth}|p{0.2\textwidth}|X|}
	\caption{Процедурные функции модуля каталога и просмотра услуг}
	\label{catalog-procedural-functions:table}\\
	\hline
	Функция & Уровень\newline доступа & Описание \\
	\hline
	\endfirsthead
	\caption*{Продолжение таблицы \ref{catalog-procedural-functions:table}}\\
	\hline
	Функция & Уровень\newline доступа & Описание \\
	\hline
	\endhead
	Получение списка услуг & public & Загружает активные услуги из базы данных с учётом сведений о категории, городе и исполнителе. \\
	\hline
	Фильтрация по тексту & public & Отбирает услуги по совпадению поисковой строки с названием или описанием услуги. \\
	\hline
	Фильтрация по категории и городу & public & Ограничивает список услуг выбранной категорией и городом оказания услуги. \\
	\hline
	Фильтрация по доступным датам & public & Отображает только те услуги, по которым имеются свободные будущие слоты в выбранные пользователем даты. \\
	\hline
	Отображение карточек услуг & public & Формирует карточки услуг с названием, стоимостью, городом, изображением и информацией о ближайшем доступном времени. \\
	\hline
	Просмотр страницы услуги & public & Отображает полную информацию об услуге, исполнителе, изображениях, отзывах и свободных слотах записи. \\
	\hline
	Отображение отзывов исполнителя & public & Загружает и показывает отзывы, связанные с исполнителем выбранной услуги. \\
	\hline
\end{xltabular}

Особенностью модуля является связь каталога с таблицей временных слотов. Благодаря этому пользователь видит не только опубликованные услуги, но и предложения, по которым действительно доступна запись на выбранные даты.

\subsubsection{Модуль управления услугами исполнителя}

Модуль управления услугами исполнителя предоставляет пользователю с ролью исполнителя возможность создавать, редактировать и просматривать собственные услуги. В системе используется запись по конкретным временным слотам, поэтому тип бронирования фиксирован и не выбирается пользователем.

Процедурные функции модуля приведены в таблице \ref{provider-services-procedural-functions:table}.

\begin{xltabular}{\textwidth}{|p{0.3\textwidth}|p{0.2\textwidth}|X|}
	\caption{Процедурные функции модуля управления услугами исполнителя}
	\label{provider-services-procedural-functions:table}\\
	\hline
	Функция & Уровень\newline доступа & Описание \\
	\hline
	\endfirsthead
	\caption*{Продолжение таблицы \ref{provider-services-procedural-functions:table}}\\
	\hline
	Функция & Уровень\newline доступа & Описание \\
	\hline
	\endhead
	Просмотр собственных услуг & provider & Отображает список услуг текущего исполнителя, их статус, изображения и переходы к редактированию или расписанию. \\
	\hline
	Отображение формы создания услуги & provider & Предоставляет исполнителю форму добавления новой услуги с выбором категории, города, цены, длительности и описания. \\
	\hline
	Проверка данных услуги & internal & Проверяет заполнение обязательных полей, корректность стоимости, длительности, категории, города и статуса. \\
	\hline
	Создание услуги & internal & Создаёт новую запись услуги в базе данных и связывает её с текущим исполнителем. \\
	\hline
	Загрузка изображений услуги & internal & Проверяет формат и размер изображений, сохраняет файлы на сервере и создаёт записи в таблице изображений услуги. \\
	\hline
	Отображение формы редактирования & provider & Загружает данные выбранной услуги и отображает форму их изменения. \\
	\hline
	Обновление услуги & internal & Проверяет права владельца услуги и сохраняет изменённые данные в базе данных. \\
	\hline
	Удаление изображений услуги & internal & Удаляет выбранные изображения из карточки услуги и обновляет связанные записи в базе данных. \\
	\hline
\end{xltabular}

В модуле предусмотрена проверка загружаемых изображений. Для одной услуги может быть загружено ограниченное количество файлов допустимого формата, что упрощает хранение и отображение материалов на странице услуги.

\subsubsection{Модуль расписания и формирования слотов}

Модуль расписания и формирования слотов предназначен для задания рабочих интервалов исполнителя и автоматического создания доступного времени для записи. Исполнитель указывает рабочие дни и временные интервалы, а система формирует свободные слоты с учётом длительности услуги.

Обычные PHP-функции модуля приведены в таблице \ref{schedule-php-functions:table}.

\begin{xltabular}{\textwidth}{|p{0.3\textwidth}|p{0.22\textwidth}|X|}
	\caption{Обычные PHP-функции модуля расписания и формирования слотов}
	\label{schedule-php-functions:table}\\
	\hline
	Функция & Область применения & Описание \\
	\hline
	\endfirsthead
	\caption*{Продолжение таблицы \ref{schedule-php-functions:table}}\\
	\hline
	Функция & Область применения & Описание \\
	\hline
	\endhead
	\texttt{generateServiceSlots(\$pdo, \$service, \$weeksValid)} & Генерация слотов & Формирует свободные временные слоты на основе расписания услуги, длительности услуги и списка недоступных дат. \\
	\hline
\end{xltabular}

Процедурные функции модуля приведены в таблице \ref{schedule-procedural-functions:table}.

\begin{xltabular}{\textwidth}{|p{0.3\textwidth}|p{0.2\textwidth}|X|}
	\caption{Процедурные функции модуля расписания и формирования слотов}
	\label{schedule-procedural-functions:table}\\
	\hline
	Функция & Уровень\newline доступа & Описание \\
	\hline
	\endfirsthead
	\caption*{Продолжение таблицы \ref{schedule-procedural-functions:table}}\\
	\hline
	Функция & Уровень\newline доступа & Описание \\
	\hline
	\endhead
	Отображение расписания услуги & provider & Показывает текущие рабочие интервалы, количество свободных и занятых слотов, а также список исключённых дат. \\
	\hline
	Проверка рабочих интервалов & internal & Проверяет корректность выбранных дней недели, времени начала, времени окончания и периода генерации. \\
	\hline
	Сохранение расписания & internal & Обновляет рабочие интервалы услуги в базе данных. \\
	\hline
	Добавление недоступной даты & provider & Добавляет дату, в которую исполнитель не принимает записи по выбранной услуге. \\
	\hline
	Удаление недоступной даты & provider & Удаляет ранее указанную исключённую дату и позволяет снова формировать слоты на этот день. \\
	\hline
	Обновление будущих слотов & internal & Удаляет будущие свободные слоты и создаёт их заново, не затрагивая уже занятые интервалы. \\
	\hline
\end{xltabular}

При генерации используется длительность услуги в минутах. Например, если рабочий интервал исполнителя задан с 10:00 до 14:00, а длительность услуги составляет 60 минут, система создаёт слоты 10:00, 11:00, 12:00 и 13:00. Такой подход позволяет автоматически поддерживать актуальный список доступного времени.

\subsubsection{Модуль бронирования услуг}

Модуль бронирования услуг отвечает за создание заявок на оказание услуги. Заказчик выбирает свободный временной слот, подтверждает запись и при необходимости оставляет комментарий. При создании бронирования система проверяет доступность выбранного слота и переводит его в состояние занятого.

Процедурные функции модуля приведены в таблице \ref{booking-procedural-functions:table}.

\begin{xltabular}{\textwidth}{|p{0.3\textwidth}|p{0.2\textwidth}|X|}
	\caption{Процедурные функции модуля бронирования услуг}
	\label{booking-procedural-functions:table}\\
	\hline
	Функция & Уровень\newline доступа & Описание \\
	\hline
	\endfirsthead
	\caption*{Продолжение таблицы \ref{booking-procedural-functions:table}}\\
	\hline
	Функция & Уровень\newline доступа & Описание \\
	\hline
	\endhead
	Отображение формы бронирования & customer & Показывает заказчику сведения об услуге, исполнителе, выбранном слоте и поле для комментария. \\
	\hline
	Проверка выбранного слота & internal & Проверяет существование слота, его принадлежность выбранной услуге, свободный статус и дату в будущем. \\
	\hline
	Проверка прав заказчика & internal & Запрещает бронирование услуги неавторизованным пользователям, исполнителям и владельцу услуги. \\
	\hline
	Создание бронирования & internal & Создаёт запись заявки в базе данных со статусом ожидания подтверждения. \\
	\hline
	Блокировка слота & internal & Блокирует выбранную строку слота на время транзакции, что снижает риск двойного бронирования. \\
	\hline
	Создание чата по заявке & internal & Создаёт чат, связанный с конкретным бронированием. \\
	\hline
	Создание системного сообщения & internal & Добавляет в чат сообщение о создании новой заявки. \\
	\hline
	Просмотр собственных бронирований & customer & Отображает список заявок текущего заказчика, их статусы и доступные действия. \\
	\hline
\end{xltabular}

Важной особенностью модуля является выполнение создания бронирования в транзакции. Выбранный слот блокируется на время проверки и сохранения данных, после чего переводится в состояние \texttt{booked}. Это предотвращает повторное бронирование одного и того же временного интервала.

\subsubsection{Модуль обработки статусов заявок}

Модуль обработки статусов заявок предназначен для изменения состояния бронирования после его создания. В системе предусмотрены статусы ожидания, подтверждения, отмены, завершения услуги и неявки заказчика.

Процедурные функции модуля приведены в таблице \ref{booking-status-procedural-functions:table}.

\begin{xltabular}{\textwidth}{|p{0.3\textwidth}|p{0.2\textwidth}|X|}
	\caption{Процедурные функции модуля обработки статусов заявок}
	\label{booking-status-procedural-functions:table}\\
	\hline
	Функция & Уровень\newline доступа & Описание \\
	\hline
	\endfirsthead
	\caption*{Продолжение таблицы \ref{booking-status-procedural-functions:table}}\\
	\hline
	Функция & Уровень\newline доступа & Описание \\
	\hline
	\endhead
	Отмена бронирования заказчиком & customer & Позволяет заказчику отменить собственную заявку, если она ещё не завершена и время записи не наступило. \\
	\hline
	Подтверждение заявки & provider & Переводит ожидающую заявку в статус подтверждённой. \\
	\hline
	Отмена заявки исполнителем & provider & Позволяет исполнителю отменить входящую заявку, если оказание услуги невозможно. \\
	\hline
	Завершение услуги & provider & Переводит подтверждённую заявку в статус завершённой после оказания услуги. \\
	\hline
	Фиксация неявки заказчика & provider & Переводит подтверждённую заявку в статус неявки, если заказчик не пришёл в назначенное время. \\
	\hline
	Освобождение слота & internal & При отмене будущей заявки переводит связанный слот обратно в свободное состояние. \\
	\hline
	Добавление системного сообщения & internal & Фиксирует изменение статуса заявки в чате бронирования. \\
	\hline
	Отказ рискованному заказчику & provider & Позволяет исполнителю отклонить заявку при наличии признаков низкой добропорядочности заказчика. \\
	\hline
\end{xltabular}

Каждое изменение статуса сопровождается системным сообщением в чате. Это позволяет заказчику и исполнителю видеть историю действий по заявке и текущее состояние бронирования.

\subsubsection{Модуль переписки по бронированию}

Модуль переписки по бронированию обеспечивает обмен сообщениями между заказчиком и исполнителем. Чат создаётся автоматически после оформления заявки и связан с конкретным бронированием.

Процедурные функции модуля приведены в таблице \ref{chat-procedural-functions:table}.

\begin{xltabular}{\textwidth}{|p{0.3\textwidth}|p{0.2\textwidth}|X|}
	\caption{Процедурные функции модуля переписки по бронированию}
	\label{chat-procedural-functions:table}\\
	\hline
	Функция & Уровень\newline доступа & Описание \\
	\hline
	\endfirsthead
	\caption*{Продолжение таблицы \ref{chat-procedural-functions:table}}\\
	\hline
	Функция & Уровень\newline доступа & Описание \\
	\hline
	\endhead
	Просмотр списка чатов & public & Отображает пользователю список переписок, связанных с его бронированиями. \\
	\hline
	Проверка доступа к чату & internal & Проверяет, что текущий пользователь является заказчиком или исполнителем по выбранному бронированию. \\
	\hline
	Просмотр истории сообщений & public & Загружает и отображает сообщения по выбранному чату. \\
	\hline
	Отправка текстового сообщения & public & Сохраняет новое текстовое сообщение пользователя в базе данных. \\
	\hline
	Загрузка изображений в чат & public & Проверяет формат и размер изображений, сохраняет их на сервере и связывает с сообщением. \\
	\hline
	Отметка сообщений как прочитанных & internal & Обновляет состояние сообщений после открытия чата пользователем. \\
	\hline
	Отображение счётчика непрочитанных сообщений & public & Показывает в интерфейсе количество непрочитанных сообщений пользователя. \\
	\hline
\end{xltabular}

Модуль исключает доступ к чужим перепискам за счёт проверки участия пользователя в бронировании. Для сообщений с изображениями используются отдельные записи, связанные с основным сообщением.

\subsubsection{Модуль избранных услуг}

Модуль избранных услуг позволяет заказчику сохранять интересующие предложения для последующего быстрого доступа. Добавление и удаление услуги выполняется одним действием, зависящим от текущего состояния записи в избранном.

Процедурные функции модуля приведены в таблице \ref{favorites-procedural-functions:table}.

\begin{xltabular}{\textwidth}{|p{0.3\textwidth}|p{0.2\textwidth}|X|}
	\caption{Процедурные функции модуля избранных услуг}
	\label{favorites-procedural-functions:table}\\
	\hline
	Функция & Уровень\newline доступа & Описание \\
	\hline
	\endfirsthead
	\caption*{Продолжение таблицы \ref{favorites-procedural-functions:table}}\\
	\hline
	Функция & Уровень\newline доступа & Описание \\
	\hline
	\endhead
	Проверка роли пользователя & internal & Разрешает работу с избранным только пользователю с ролью заказчика. \\
	\hline
	Проверка состояния услуги & internal & Проверяет существование и активность услуги перед добавлением в избранное. \\
	\hline
	Добавление услуги в избранное & customer & Создаёт связь между заказчиком и выбранной услугой. \\
	\hline
	Удаление услуги из избранного & customer & Удаляет ранее созданную связь между заказчиком и услугой. \\
	\hline
	Просмотр избранных услуг & customer & Отображает список сохранённых услуг с основной информацией и переходом к карточке услуги. \\
	\hline
\end{xltabular}

\subsubsection{Модуль отзывов и репутации}

Модуль отзывов и репутации предназначен для оценки взаимодействия между пользователями после завершения услуги. Заказчик может оставить отзыв об исполнителе, а исполнитель -- отзыв о заказчике после завершения услуги или фиксации неявки.

Процедурные функции модуля приведены в таблице \ref{reviews-procedural-functions:table}.

\begin{xltabular}{\textwidth}{|p{0.3\textwidth}|p{0.2\textwidth}|X|}
	\caption{Процедурные функции модуля отзывов и репутации}
	\label{reviews-procedural-functions:table}\\
	\hline
	Функция & Уровень\newline доступа & Описание \\
	\hline
	\endfirsthead
	\caption*{Продолжение таблицы \ref{reviews-procedural-functions:table}}\\
	\hline
	Функция & Уровень\newline доступа & Описание \\
	\hline
	\endhead
	Отображение формы отзыва & public & Показывает форму оценки и комментария по конкретному бронированию. \\
	\hline
	Определение получателя отзыва & internal & Определяет, какого пользователя может оценить текущий участник бронирования. \\
	\hline
	Проверка права на отзыв & internal & Проверяет участие пользователя в бронировании и допустимый статус заявки. \\
	\hline
	Проверка повторного отзыва & internal & Запрещает одному пользователю оставлять несколько отзывов по одному бронированию. \\
	\hline
	Сохранение отзыва & internal & Сохраняет оценку, текст комментария, автора и получателя отзыва в базе данных. \\
	\hline
	Формирование рейтинга исполнителя & public & Использует накопленные отзывы для отображения средней оценки исполнителя в каталоге и на странице услуги. \\
	\hline
	Оценка добропорядочности заказчика & internal & Использует отзывы исполнителей для определения репутации заказчика. \\
	\hline
\end{xltabular}

\subsubsection{Модуль оформления интерфейса}

Модуль оформления интерфейса отвечает за визуальное представление страниц web-сервиса. В него входят общие стили, стили отдельных разделов и JavaScript-код, связанный с пользовательским интерфейсом.

Процедурные функции модуля приведены в таблице \ref{ui-procedural-functions:table}.

\begin{xltabular}{\textwidth}{|p{0.3\textwidth}|p{0.2\textwidth}|X|}
	\caption{Процедурные функции модуля оформления интерфейса}
	\label{ui-procedural-functions:table}\\
	\hline
	Функция & Уровень\newline доступа & Описание \\
	\hline
	\endfirsthead
	\caption*{Продолжение таблицы \ref{ui-procedural-functions:table}}\\
	\hline
	Функция & Уровень\newline доступа & Описание \\
	\hline
	\endhead
	Формирование общего оформления & public & Задаёт базовые цвета, отступы, контейнеры, кнопки, карточки и типовые элементы интерфейса. \\
	\hline
	Оформление страниц авторизации & public & Определяет внешний вид форм регистрации и входа. \\
	\hline
	Оформление каталога услуг & public & Задаёт стили карточек услуг, фильтров, страницы услуги и форм управления услугами. \\
	\hline
	Оформление расписания & public & Определяет внешний вид страницы настройки рабочих интервалов и исключённых дат. \\
	\hline
	Оформление бронирований & public & Задаёт внешний вид страниц создания и просмотра заявок. \\
	\hline
	Оформление чата & public & Определяет внешний вид списка чатов, сообщений и формы отправки. \\
	\hline
	Оформление отзывов & public & Задаёт стили формы отзыва и блоков отображения оценок. \\
	\hline
	Переключение темы интерфейса & public & Позволяет пользователю переключать светлую и тёмную тему с сохранением выбранного режима в браузере. \\
	\hline
\end{xltabular}

\subsection{Системное тестирование}

Системное тестирование проводилось в соответствии со сценариями использования, определёнными в техническом задании. Проверялись основные функции web-сервиса: регистрация и авторизация пользователей, управление профилем, создание услуги, настройка расписания, формирование слотов, поиск услуг, бронирование, изменение статусов заявок, переписка, избранное, отзывы и обработка ошибочных ситуаций.

Тестирование выполнялось в локальной среде Open Server Panel с использованием web-сервера Apache, интерпретатора PHP и базы данных MySQL. Для проверки интерфейса использовался web-браузер. Сводная таблица тест-кейсов приведена в таблице \ref{test-cases-summary:table}.

\subsubsection{Сводная таблица тест-кейсов}

\begin{xltabular}{\textwidth}{|p{0.12\textwidth}|p{0.58\textwidth}|X|}
	\caption{Сводная таблица тест-кейсов}
	\label{test-cases-summary:table}\\
	\hline
	ID & Название тест-кейса & Результат \\
	\hline
	\endfirsthead
	\caption*{Продолжение таблицы \ref{test-cases-summary:table}}\\
	\hline
	ID & Название тест-кейса & Результат \\
	\hline
	\endhead
	TC-01 & Регистрация пользователя с ролью заказчика & Пройден \\
	\hline
	TC-02 & Авторизация зарегистрированного пользователя & Пройден \\
	\hline
	TC-03 & Редактирование профиля пользователя & Пройден \\
	\hline
	TC-04 & Создание услуги исполнителем & Пройден \\
	\hline
	TC-05 & Настройка расписания исполнителя & Пройден \\
	\hline
	TC-06 & Формирование свободных временных слотов & Пройден \\
	\hline
	TC-07 & Поиск и фильтрация услуг в каталоге & Пройден \\
	\hline
	TC-08 & Создание бронирования заказчиком & Пройден \\
	\hline
	TC-09 & Защита от повторного бронирования занятого слота & Пройден \\
	\hline
	TC-10 & Подтверждение заявки исполнителем & Пройден \\
	\hline
	TC-11 & Отмена бронирования заказчиком & Пройден \\
	\hline
	TC-12 & Переписка по заявке & Пройден \\
	\hline
	TC-13 & Добавление услуги в избранное & Пройден \\
	\hline
	TC-14 & Оставление отзыва после завершения услуги & Пройден \\
	\hline
	TC-15 & Проверка валидации некорректных данных & Пройден \\
	\hline
\end{xltabular}

\subsubsection{TC-01. Регистрация пользователя с ролью заказчика}

Предусловие: пользователь не авторизован в системе; открыта страница регистрации.

Шаги:
\begin{enumerate}
	\item Открыть страницу регистрации.
	\item Заполнить ФИО, email, телефон, город и пароль.
	\item Выбрать роль заказчика.
	\item Подтвердить пароль.
	\item Нажать кнопку регистрации.
\end{enumerate}

Ожидаемый результат: система проверяет корректность введённых данных, создаёт новую учётную запись и перенаправляет пользователя на страницу входа.

Фактический результат: учётная запись была создана, пользователь получил сообщение об успешной регистрации и был перенаправлен на страницу авторизации.

На рисунке~\ref{tc01-register:image} представлен результат регистрации пользователя с ролью заказчика.

\begin{figure}[H]
	\center{\includegraphics[width=1\linewidth]{tc1.png}}
	\caption{Результат регистрации пользователя}
	\label{tc01-register:image}
\end{figure}

\subsubsection{TC-02. Авторизация зарегистрированного пользователя}

Предусловие: в системе существует зарегистрированный пользователь.

Шаги:
\begin{enumerate}
	\item Открыть страницу входа.
	\item Ввести email и пароль зарегистрированного пользователя.
	\item Нажать кнопку входа.
\end{enumerate}

Ожидаемый результат: система находит пользователя в базе данных, проверяет пароль и создаёт пользовательскую сессию.

Фактический результат: пользователь успешно авторизован и перенаправлен в личный кабинет.

На рисунке~\ref{tc02-login:image} представлен результат авторизации зарегистрированного пользователя.

\begin{figure}[H]
	\center{\includegraphics[width=1\linewidth]{tc2.png}}
	\caption{Результат авторизации пользователя}
	\label{tc02-login:image}
\end{figure}

\subsubsection{TC-03. Редактирование профиля пользователя}

Предусловие: пользователь авторизован в системе.

Шаги:
\begin{enumerate}
	\item Открыть страницу редактирования профиля.
	\item Изменить телефон, город или описание пользователя.
	\item При необходимости выбрать изображение профиля.
	\item Сохранить изменения.
\end{enumerate}

Ожидаемый результат: система проверяет введённые данные, сохраняет изменения и отображает обновлённые сведения в профиле.

Фактический результат: данные профиля были успешно обновлены, пользователь получил сообщение об успешном сохранении.

На рисунке~\ref{tc03-profile:image} представлен результат редактирования профиля пользователя.

\begin{figure}[H]
	\center{\includegraphics[width=1\linewidth]{tc3.png}}
	\caption{Результат редактирования профиля пользователя}
	\label{tc03-profile:image}
\end{figure}

\subsubsection{TC-04. Создание услуги исполнителем}

Предусловие: пользователь авторизован с ролью исполнителя.

Шаги:
\begin{enumerate}
	\item Открыть раздел управления услугами исполнителя.
	\item Перейти к форме добавления услуги.
	\item Заполнить название услуги, категорию, город, цену, длительность, статус, адрес и описание.
	\item Нажать кнопку сохранения услуги.
\end{enumerate}

Ожидаемый результат: система проверяет корректность данных, создаёт карточку услуги и отображает её в списке услуг исполнителя.

Фактический результат: услуга была успешно создана и появилась в списке услуг исполнителя.

На рисунке~\ref{tc04-add-service:image} представлен результат создания услуги исполнителем.

\begin{figure}[H]
	\center{\includegraphics[width=1\linewidth]{tc4.png}}
	\caption{Создание услуги исполнителем}
	\label{tc04-add-service:image}
\end{figure}

\subsubsection{TC-05. Настройка расписания исполнителя}

Предусловие: пользователь авторизован с ролью исполнителя; у исполнителя создана услуга.

Шаги:
\begin{enumerate}
	\item Открыть страницу настройки расписания выбранной услуги.
	\item Выбрать рабочие дни недели.
	\item Указать время начала и окончания работы.
	\item Указать количество недель, на которое необходимо сформировать расписание.
	\item Сохранить расписание.
\end{enumerate}

Ожидаемый результат: система сохраняет рабочие интервалы и запускает формирование свободных временных слотов.

Фактический результат: расписание было сохранено, система сформировала доступные слоты для записи.

На рисунке~\ref{tc05-schedule:image} представлен результат настройки расписания исполнителя.

\begin{figure}[H]
	\center{\includegraphics[width=1\linewidth]{tc5.png}}
	\caption{Настройка расписания исполнителя}
	\label{tc05-schedule:image}
\end{figure}

\subsubsection{TC-06. Формирование свободных временных слотов}

Предусловие: для услуги задано расписание и указана длительность услуги.

Шаги:
\begin{enumerate}
	\item Сохранить расписание услуги.
	\item Открыть страницу услуги в каталоге.
	\item Проверить список доступных дат и времени записи.
\end{enumerate}

Ожидаемый результат: система отображает свободные слоты, сформированные на основе расписания и длительности услуги.

Фактический результат: на странице услуги были отображены доступные временные слоты для бронирования.

На рисунке~\ref{tc06-slots:image} представлен результат формирования свободных временных слотов.

\begin{figure}[H]
	\center{\includegraphics[width=1\linewidth]{tc6.png}}
	\caption{Сформированные временные слоты}
	\label{tc06-slots:image}
\end{figure}

\subsubsection{TC-07. Поиск и фильтрация услуг в каталоге}

Предусловие: в системе имеются активные услуги со свободными слотами.

Шаги:
\begin{enumerate}
	\item Открыть каталог услуг.
	\item Ввести поисковую строку или выбрать категорию.
	\item Выбрать город.
	\item Выбрать одну или несколько доступных дат.
	\item Применить фильтры.
\end{enumerate}

Ожидаемый результат: система отображает только те услуги, которые соответствуют выбранным параметрам и имеют свободные слоты на выбранные даты.

Фактический результат: каталог был отфильтрован, пользователю отображены подходящие услуги с информацией о ближайшем доступном времени.

На рисунке~\ref{tc07-filter:image} представлен результат поиска и фильтрации услуг в каталоге.

\begin{figure}[H]
	\center{\includegraphics[width=1\linewidth]{tc7.png}}
	\caption{Фильтрация услуг в каталоге}
	\label{tc07-filter:image}
\end{figure}

\subsubsection{TC-08. Создание бронирования заказчиком}

Предусловие: пользователь авторизован с ролью заказчика; на странице услуги имеется свободный слот.

Шаги:
\begin{enumerate}
	\item Открыть страницу услуги.
	\item Выбрать свободный временной слот.
	\item Перейти на страницу подтверждения бронирования.
	\item При необходимости ввести комментарий заказчика.
	\item Подтвердить создание заявки.
\end{enumerate}

Ожидаемый результат: система создаёт бронирование со статусом ожидания, закрепляет выбранный слот за заказчиком и создаёт чат по заявке.

Фактический результат: заявка была успешно создана, слот получил статус занятого, по бронированию был создан чат.

На рисунке~\ref{tc08-booking:image} представлен результат создания бронирования заказчиком.

\begin{figure}[H]
	\center{\includegraphics[width=1\linewidth]{tc8.png}}
	\caption{Создание бронирования заказчиком}
	\label{tc08-booking:image}
\end{figure}

\subsubsection{TC-09. Защита от повторного бронирования занятого слота}

Предусловие: в системе существует слот, который уже был забронирован другим пользователем.

Шаги:
\begin{enumerate}
	\item Попытаться повторно открыть бронирование уже занятого слота.
	\item Отправить форму создания бронирования.
\end{enumerate}

Ожидаемый результат: система не создаёт повторную заявку и выводит сообщение о том, что слот уже занят или недоступен.

Фактический результат: повторное бронирование не было создано, пользователь получил сообщение об ошибке.

На рисунке~\ref{tc9-double-booking:image} представлен результат проверки защиты от повторного бронирования занятого слота.

\begin{figure}[H]
	\center{\includegraphics[width=1\linewidth]{tc9.png}}
	\caption{Проверка повторного бронирования занятого слота}
	\label{tc9-double-booking:image}
\end{figure}

\subsubsection{TC-10. Подтверждение заявки исполнителем}

Предусловие: пользователь авторизован с ролью исполнителя; по его услуге существует заявка в статусе ожидания.

Шаги:
\begin{enumerate}
	\item Открыть раздел входящих заявок исполнителя.
	\item Найти заявку в статусе ожидания.
	\item Нажать кнопку подтверждения заявки.
\end{enumerate}

Ожидаемый результат: система переводит заявку в статус подтверждённой и добавляет системное сообщение в чат бронирования.

Фактический результат: статус заявки был изменён на подтверждённый, в чате появилось системное сообщение о подтверждении.

На рисунке~\ref{tc10-confirm:image} представлен результат подтверждения заявки исполнителем.

\begin{figure}[H]
	\center{\includegraphics[width=1\linewidth]{tc10.png}}
	\caption{Подтверждение заявки исполнителем}
	\label{tc10-confirm:image}
\end{figure}

\subsubsection{TC-11. Отмена бронирования заказчиком}

Предусловие: пользователь авторизован с ролью заказчика; у пользователя есть заявка в статусе ожидания или подтверждения, время которой ещё не наступило.

Шаги:
\begin{enumerate}
	\item Открыть раздел своих бронирований.
	\item Выбрать активную заявку.
	\item Нажать кнопку отмены бронирования.
	\item Подтвердить действие.
\end{enumerate}

Ожидаемый результат: система переводит заявку в статус отменённой, фиксирует инициатора отмены и освобождает будущий слот.

Фактический результат: заявка была отменена, слот снова стал доступен для бронирования.

На рисунке~\ref{tc11-cancel:image} представлен результат отмены бронирования заказчиком.

\begin{figure}[H]
	\center{\includegraphics[width=1\linewidth]{tc11.png}}
	\caption{Отмена бронирования заказчиком}
	\label{tc11-cancel:image}
\end{figure}

\subsubsection{TC-12. Переписка по заявке}

Предусловие: пользователь является участником бронирования; по заявке создан чат.

Шаги:
\begin{enumerate}
	\item Открыть список чатов.
	\item Перейти в чат по выбранному бронированию.
	\item Ввести текст сообщения.
	\item При необходимости прикрепить изображение.
	\item Отправить сообщение.
\end{enumerate}

Ожидаемый результат: система сохраняет сообщение, отображает его в истории переписки и прикрепляет выбранные изображения.

Фактический результат: сообщение было успешно отправлено и отображено в чате по бронированию.

На рисунке~\ref{tc12-chat:image} представлен результат переписки по заявке.

\begin{figure}[H]
	\center{\includegraphics[width=1\linewidth]{tc12.png}}
	\caption{Переписка по заявке}
	\label{tc12-chat:image}
\end{figure}

\subsubsection{TC-13. Добавление услуги в избранное}

Предусловие: пользователь авторизован с ролью заказчика; открыта карточка услуги или страница услуги.

Шаги:
\begin{enumerate}
	\item Нажать кнопку добавления услуги в избранное.
	\item Открыть раздел избранных услуг.
	\item Проверить наличие выбранной услуги в списке.
\end{enumerate}

Ожидаемый результат: система сохраняет связь между пользователем и услугой, после чего услуга отображается в разделе избранного.

Фактический результат: услуга была добавлена в избранное и отображена в соответствующем разделе.

На рисунке~\ref{tc13-favorites:image} представлен результат добавления услуги в избранное.

\begin{figure}[H]
	\center{\includegraphics[width=1\linewidth]{tc13.png}}
	\caption{Добавление услуги в избранное}
	\label{tc13-favorites:image}
\end{figure}

\subsubsection{TC-14. Оставление отзыва после завершения услуги}

Предусловие: бронирование переведено в статус завершённой услуги; пользователь является участником данного бронирования.

Шаги:
\begin{enumerate}
	\item Открыть завершённое бронирование.
	\item Перейти к форме создания отзыва.
	\item Указать оценку и текст комментария.
	\item Отправить отзыв.
\end{enumerate}

Ожидаемый результат: система проверяет право пользователя на создание отзыва, сохраняет отзыв и обновляет репутационную информацию пользователя, которому он адресован.

Фактический результат: отзыв был успешно сохранён и стал учитываться при отображении рейтинга пользователя.

На рисунке~\ref{tc14-review:image} представлен результат оставления отзыва после завершения услуги.

\begin{figure}[H]
	\center{\includegraphics[width=1\linewidth]{tc14.png}}
	\caption{Оставление отзыва после завершения услуги}
	\label{tc14-review:image}
\end{figure}

\subsubsection{TC-15. Проверка валидации некорректных данных}

Предусловие: пользователь открыл одну из форм системы, например форму регистрации, создания услуги или отправки сообщения.

Шаги:
\begin{enumerate}
	\item Оставить обязательные поля пустыми или ввести некорректные значения.
	\item Отправить форму.
	\item Проверить реакцию системы.
\end{enumerate}

Ожидаемый результат: система не сохраняет некорректные данные, выводит сообщение об ошибке и оставляет пользователя на текущей странице для исправления данных.

Фактический результат: некорректные данные не были сохранены, пользователь получил сообщения об ошибках.

На рисунке~\ref{tc15-validation:image} представлен результат проверки валидации некорректных данных.

\begin{figure}[H]
	\center{\includegraphics[width=1\linewidth]{tc15.png}}
	\caption{Проверка валидации некорректных данных}
	\label{tc15-validation:image}
\end{figure}

\subsubsection{Нефункциональное тестирование}

Помимо проверки основных пользовательских сценариев было выполнено нефункциональное тестирование. Проверялась корректность отображения страниц в браузере, обработка ошибочных действий пользователя, защита закрытых разделов от неавторизованного доступа, разграничение функций заказчика и исполнителя, а также корректность обработки загружаемых файлов.

Результаты нефункционального тестирования приведены в таблице \ref{non-functional-tests:table}.

\begin{xltabular}{\textwidth}{|p{0.38\textwidth}|p{0.42\textwidth}|X|}
	\caption{Результаты нефункционального тестирования}
	\label{non-functional-tests:table}\\
	\hline
	Проверка & Описание & Результат \\
	\hline
	\endfirsthead
	\caption*{Продолжение таблицы \ref{non-functional-tests:table}}\\
	\hline
	Проверка & Описание & Результат \\
	\hline
	\endhead
	Проверка доступа к закрытым страницам & Неавторизованный пользователь перенаправляется на страницу входа при попытке открыть личный кабинет, бронирования, чат или управление услугами. & Пройден \\
	\hline
	Разграничение ролей & Заказчик не может управлять услугами исполнителя, а исполнитель не может бронировать собственные услуги. & Пройден \\
	\hline
	Безопасный вывод данных & Пользовательские данные выводятся через функцию экранирования, что снижает риск выполнения HTML- или JavaScript-кода. & Пройден \\
	\hline
	Проверка загружаемых файлов & Система принимает только допустимые изображения и отклоняет файлы неподходящего формата или размера. & Пройден \\
	\hline
	Сохранение целостности данных & При ошибках во время создания бронирования или загрузки файлов операция завершается корректно, а данные не сохраняются частично. & Пройден \\
	\hline
\end{xltabular}

\subsubsection{Вывод по результатам тестирования}

В ходе системного тестирования были проверены основные сценарии работы web-сервиса для размещения объявлений об услугах. Все тест-кейсы завершились успешно. Система корректно выполняет регистрацию и авторизацию пользователей, управление профилями, публикацию услуг, настройку расписания, формирование свободных слотов, поиск и фильтрацию услуг, создание бронирований, обработку статусов заявок, переписку, работу с избранным и отзывами.

Проверка ошибочных ситуаций показала, что система выполняет валидацию пользовательских данных, ограничивает доступ к защищённым разделам и предотвращает повторное бронирование занятого временного слота. Таким образом, разработанная программная система соответствует функциональным требованиям, сформулированным в техническом задании.